Wir zeigen, dass es zu jedem Punkt

\mavergleichskette
{\vergleichskette
{ x
}
{ \in }{ U
}
{ = }{ X \setminus Y
}
{    }{
}
{   }{
}
}
{}{}{}
eine offene Umgebung

\mavergleichskette
{\vergleichskette
{ x
}
{ \in }{ V
}
{ \subseteq }{ U
}
{    }{
}
{   }{
}
}
{}{}{}
gibt. Deren Vereinigung ist dann gleich $U$. Sei also

\mavergleichskette
{\vergleichskette
{ x
}
{ \notin }{ Y
}
{  }{
}
{    }{
}
{   }{
}
}
{}{}{}
gegeben. Zu jedem Punkt

\mavergleichskette
{\vergleichskette
{ y
}
{ \in }{ Y
}
{  }{
}
{    }{
}
{   }{
}
}
{}{}{}
gibt es offene Mengen

\mavergleichskette
{\vergleichskette
{ x
}
{ \in }{ V_y
}
{  }{
}
{    }{
}
{   }{
}
}
{}{}{}
und

\mavergleichskette
{\vergleichskette
{ y
}
{ \in }{ W_y
}
{  }{
}
{    }{
}
{   }{
}
}
{}{}{,}
die zueinander disjunkt sind. Die offenen Mengen

\mathbed {W_y} {}
 {y \in Y} {}
 {} {} {} {,}
überdecken $Y$. Aufgrund der Kompaktheit gibt es davon eine endliche Teilüberdeckung, sagen wir

\mavergleichskettedisp
{\vergleichskette
{ Y
}
{ \subseteq} { \bigcup_{i = 1}^n  W_i
}
{ } {
}
{ } {
}
{ } {
}
}
{}{}{.}
Dann ist

\mavergleichskette
{\vergleichskette
{ V
}
{ = }{ \bigcap_{i = 1}^n  V_i
}
{  }{
}
{    }{
}
{   }{
}
}
{}{}{}
als endlicher Durchschnitt von offenen Mengen offen und damit eine offene Umgebung von $x$, die zu $Y$ disjunkt ist.

 [[Kategorie:Latexseite]]
